% GERADO por tools/generation/gerar_secoes_latex.py a partir de
% artifacts/paper/secoes/03-trabalhos-relacionados.md. Nao editar a mao: a proxima geracao
% desfaz. Corrija o Markdown e regere.
Ontologias já foram empregadas em contextos de observatório, sempre como meio de integrar e recuperar os dados do fenômeno observado. Numa delas, o casamento semântico entre pedidos de usuários e serviços de processamento de imagens sustenta um observatório virtual multiagente \cite{chen2008apply}; o \textit{Virtual Solar-Terrestrial Observatory} organiza dados de física solar sob ontologias que ocultam a heterogeneidade das fontes \cite{fox2009ontology}; no monitoramento ambiental, uma ontologia modular fundamentada na Basic Formal Ontology integra desastres, sensores e medições e habilita inferência sobre dados de precipitação \cite{masmoudi2018ontology}; e, na saúde, um modelo conceitual de observatórios incorpora camada semântica para estruturar o conhecimento que sustenta políticas públicas \cite{yoshiura2018towards}. Em todos, a ontologia é infraestrutura para os dados observados; nenhum toma como objeto o modelo conceitual que descreve o próprio observatório.

Na direção complementar, há trabalhos que tomam modelos conceituais como objeto e os avaliam ontologicamente. A Representation Theory foi aplicada a gramáticas de modelagem de processos para identificar deficiências representacionais e medir seu efeito percebido sobre quem as usa \cite{recker2011ontological}. Na linha da UFO, a avaliação ontológica de linguagens e modelos é o método que fundamenta a própria OntoUML \cite{guizzardi2005ontological}, e há ontologias de referência construídas por reconciliação de um modelo de domínio preexistente com as categorias da UFO --- na cardiologia \cite{gonccalves2011using}, no direito penal \cite{mario2020handling} e na administração pública brasileira, em que a autorização orçamentária e a execução da despesa recebem ontologia de referência publicada na iSys \cite{detoni2019ontologia}. Esses trabalhos demonstram o método; nenhum deles tem por objeto um modelo de referência do domínio de observatórios.

A lacuna que este artigo ocupa está no cruzamento das duas linhas, e tem três marcas. Primeira: nenhum trabalho submete um modelo de referência publicado do domínio de observatórios a uma análise ontológica crítica. Segunda: nos trabalhos que analisam modelos, o resultado é a ontologia construída; aqui o resultado declarado são as deficiências que a formalização publicada expôs, classificadas por uma teoria de SI, com cada achado rastreado até o ponto do texto do modelo de referência que o permitiu --- é o que converte correção de um artefato em conhecimento sobre o modelo que o originou. Terceira: a evidência é a diferença entre um estado anterior e um posterior, ambos publicados e reexecutáveis, e não uma afirmação sobre a qualidade do resultado final.
